% Generated by tools/annotation/descriptives.py from the corpus manifest, the split
% registry, the frozen A2 frame manifest and the registered A2 sample manifest. Do not edit by hand.
% Appendix A.2 since 2026-11-09 (slimming debate consensus F-03): the main text carries the summary tab:corpus.
\begin{table}[htbp]
\centering
\footnotesize
\setlength{\tabcolsep}{3pt}
\begin{threeparttable}
\caption[Corpus, eligible frame and evaluation-sample allocation by outlet]{Corpus, eligible frame and evaluation-sample allocation by outlet; each outlet with eligible paragraphs is one stratum within its language, with stratum size $N_h$ in the eligible frame, draw $n_h$ and design weight $w_h = N_h/n_h$.}
\label{tab:corpus_outlets}
\begin{tabular}{@{}llrrr@{\hspace{12pt}}rrr@{}}
\toprule
 & & \multicolumn{3}{c}{\textbf{Corpus}} & \multicolumn{3}{c}{\textbf{Stratum}} \\
\cmidrule(lr){3-5}\cmidrule(l){6-8}
\textbf{Outlet} & \textbf{Country} & \textbf{Articles} & \textbf{Paragraphs} & \textbf{Candidates} & $\boldsymbol{N_h}$ & $\boldsymbol{n_h}$ & $\boldsymbol{w_h}$ \\
\midrule
\multicolumn{8}{@{}l}{\textit{German-language outlets}} \\
taz & Germany & 8 & 80 & 64 & 45 & 12 & 3.75 \\
DIE ZEIT & Germany & 7 & 52 & 44 & 24 & 7 & 3.43 \\
Süddeutsche Zeitung & Germany & 4 & 40 & 23 & 23 & 6 & 3.83 \\
BILD & Germany & 6 & 38 & 28 & 28 & 8 & 3.50 \\
Neue Zürcher Zeitung & Switzerland & 2 & 36 & 22 & 22 & 6 & 3.67 \\
tagesschau.de & Germany & 3 & 31 & 24 & 24 & 6 & 4.00 \\
Handelsblatt & Germany & 2 & 26 & 12 & 12 & 3 & 4.00 \\
DER SPIEGEL & Germany & 3 & 21 & 18 & 18 & 5 & 3.60 \\
Hamburger Abendblatt & Germany & 1 & 11 & 9 & 9 & 2 & 4.50 \\
BR24 & Germany & 1 & 8 & 7 & 7 & 2 & 3.50 \\
Berliner Morgenpost & Germany & 1 & 7 & 4 & 4 & 1 & 4.00 \\
STERN.de & Germany & 1 & 7 & 6 & 5 & 1 & 5.00 \\
Salzburger Nachrichten (SN.at) & Austria & 1 & 4 & 4 & 4 & 1 & 4.00 \\
\textit{All German} & & 40 & 361 & 265 & 225 & 60 & \\
\addlinespace
\multicolumn{8}{@{}l}{\textit{English-language outlets}} \\
The New York Times & United States & 13 & 311 & 184 & 184 & 24 & 7.67 \\
The Economist & United Kingdom & 12 & 162 & 125 & 108 & 18 & 6.00 \\
The Guardian & United Kingdom & 4 & 95 & 68 & 68 & 12 & 5.67 \\
Al Jazeera & Qatar & 3 & 94 & 37 & 22 & 4 & 5.50 \\
The Washington Post & United States & 1 & 18 & 11 & 11 & 2 & 5.50 \\
Clean Energy Wire & Germany & 1 & 6 & 6 & 0 & -- & -- \\
\textit{All English} & & 34 & 686 & 431 & 393 & 60 & \\
\midrule
Total & 6 countries & 74 & 1{,}047 & 696 & 618 & 120 & \\
\bottomrule
\end{tabular}
\begin{tablenotes}\footnotesize
\item Articles: articles with body prose (articles without it are not counted). Candidates: paragraphs the keyword prefilter marks as naming at least one studied technology. $N_h$: candidates in the reserve that carry no transcription gap and do not come from the article quoted in the prompt. Clean Energy Wire forms no stratum: its one article is the one quoted in the prompt. Allocation within language by the largest-remainder method, cap 24 per stratum (binding only where $n_h = 24$).
\end{tablenotes}
\end{threeparttable}
\end{table}
